% ═══════════════════════════════════════════════════════════════════════════════
% CASO 2: REFINERÍA BICENTENARIO - INSTITUTO ATENEO DE COLIMA
% Estudio Documental: Corrupción en México (2006-2012)
% ═══════════════════════════════════════════════════════════════════════════════

\nota{
  Este capítulo analiza los patrones de vulnerabilidad en megainfraestructura 
  del caso Refinería Bicentenario mediante evidencia de auditorías ASF (2015, 2016) 
  y doctrina de Cancino/BID (2020) y Sandoval Ballesteros (2024).
}

\section{Anuncio político, competencia estatal y endeudamiento subnacional (2008--2012)}

El 18 de marzo de 2008, durante la conmemoración de la expropiación petrolera, el Ejecutivo federal anunció la edificación de una nueva refinería con una inversión estimada superior a 9,000 millones de dólares. Para seleccionar la sede, Petróleos Mexicanos (Pemex) instrumentó un proceso competitivo entre entidades federativas ---en el que Hidalgo y Guanajuato llegaron como finalistas--- cuyo criterio decisivo se concentró en la donación estatal de un predio de 700 hectáreas con infraestructura básica de acceso.

Con el propósito de obtener la adjudicación de la sede en Tula de Allende, el gobierno del estado de Hidalgo adquirió parcelas ejidales y financió obras de conectividad mediante la contratación de deuda pública cercana a 1,500 millones de pesos (1,050 millones para la adquisición de tierras y 450 millones para infraestructura de acceso). Dicho compromiso financiero gravó las finanzas locales durante más de una década: la administración estatal liquidó por completo el adeudo con Banorte hasta el 18 de septiembre de 2025, diecisiete años después del anuncio presidencial y sin que la planta refinadora hubiera sido construida \citep{periodicohidalgo2025}.

Simultáneamente, entre 2009 y 2012, Pemex ejerció más de 9,000 millones de pesos (más de 620 millones de dólares) en estudios de ingeniería básica, limpieza, nivelación del terreno y la edificación de una barda perimetral con costo superior a 500 millones de pesos, sin consolidar el esquema financiero definitivo ni iniciar la construcción de las plantas industriales de proceso \citep{asf2015,asf2016,raziel2021}.

\section{Cancelación, fiscalización de la ASF y pasivos persistentes (2014--2025)}

En diciembre de 2014, durante la administración de Enrique Peña Nieto, Pemex canceló definitivamente la Refinería Bicentenario al argumentar la ausencia de rentabilidad económica frente a la coyuntura energética internacional \citep{montoya2014}. La fiscalización realizada por la Auditoría Superior de la Federación documentó deficiencias sustantivas en dos auditorías específicas:

\begin{itemize}
  \item \textbf{Auditoría de Inversiones Físicas 13-1-18T4M-04-0353 (DE-226)} sobre el Estudio de Preinversión \citep{asf2015}: Documentó inconsistencias técnicas y observaciones de fiscalización desde la etapa temprana de evaluación de factibilidad económica y de ingeniería del proyecto.
  
  \item \textbf{Auditoría 533-DE} sobre el terreno donado en Tula de Allende \citep{asf2016}: Verificó la condición improductiva de las 700 hectáreas donadas por Hidalgo y las omisiones en el resguardo patrimonial de la inversión ejercida.
\end{itemize}

Los costos fiscales continuaron acumulándose tras la cancelación: reportes sustentados en estados financieros auditados evidenciaron que Pemex siguió cubriendo pasivos vinculados al predio sin uso productivo, erogando 195.5 millones de pesos únicamente en el ejercicio 2020 \citep{raziel2021}. Como postura de descargo, en febrero de 2025 Felipe Calderón sostuvo que su administración cumplió la etapa inicial programada (ingeniería, preparación del terreno y barda perimetral) y atribuyó la pérdida del proyecto a la orden de cancelación emitida por su sucesor \citep{elfinanciero2025}.

\section{Panorama comparado de otros expedientes de obra y contratación del sexenio (2006--2012)}

La triangulación documental mostró que las vulnerabilidades observadas en Tula coexistieron con otros expedientes relevantes de obra pública, asociaciones público-privadas y sobornos corporativos registrados durante el mismo periodo:

\begin{table}[H]
\centering
\caption{Panorama comparado de otros expedientes de obra y contratación del sexenio (2006--2012)}
\label{tab:expedientes}
\small
\begin{tabular}{@{}p{2.5cm}p{4cm}p{6.5cm}@{}}
\toprule
\textbf{Expediente (2006--2012)} & \textbf{Tipología / mecanismo} & \textbf{Evidencia oficial y magnitud documentada} \\
\midrule
Estela de Luz & Sobreprecios y pagos improcedentes en obra pública. & La ASF (2013) acreditó un sobrecosto de 192\,\% (de \textasciitilde393--398 mdp a >1,104--1,146 mdp), 15 meses de retraso y 399.2 mdp en pagos improcedentes de acero \citep{asf2013,obrasexpansion2013}. \\
\addlinespace
Contratos penitenciarios (CPS) & APP opacas con pagos rígidos a 20--22 años. & Ocho CEFERESOS desde 2010; la ASF observó falta de fundamento legal claro e informes documentaron cláusulas de secrecía y pagos independientes de la ocupación real \citep{documenta2022,badillo2021}. \\
\addlinespace
Sobornos FCPA (ABB/CFE y BizJet) & Cohecho transnacional sancionado por autoridades de EE.\,UU. & Multas de 58.8 MDD (ABB/CFE) y 11.8 MDD (BizJet) por sobornos a mandos de Policía Federal y flota presidencial, con declaración de culpabilidad adicional en octubre de 2026 \citep{doj2013bizjet,doj2026bizjet,rodriguez2011}. \\
\addlinespace
Pemex: Odebrecht y Oceanografía & Continuidad transexenal y tráfico de influencias. & Contratos en Minatitlán y Etileno XXI (Odebrecht) y 54 contratos con Oceanografía entre 2002 y 2012, cuyas tramas se extendieron y judicializaron en sexenios posteriores \citep{fundar2022,guerrero2018}. \\
\bottomrule
\end{tabular}
\end{table}